\section{Conclusion}
\label{sec:conclusion}

Courts that compel governments to buy medications save individual lives.
They also raise public procurement costs by \BPnegPctHeadline{} per
purchase across the full sample, attract \BPfirmsPctHeadline{} fewer
bidders, and---once selection into the SES/SP administrative channel is
bounded---generate an under-the-gun gap of $[\BPutgBoundLow, \BPutgBoundHigh]$
between litigated and administrative urgent purchases of the same item.

The empirical signature of accountability-induced inefficiency in our
data is \emph{sourcing}, not \emph{pricing}. Within firm-buyer-item
triples observed under both the litigated and the administrative
regime, the same firm sells the same item to the same buyer at
essentially the same per-unit price ($\BPutgWithinFirmOffset$, NULL).
Suppliers that already serve the buyer's urgent procurement do not
extract a markup when the official is sanction-exposed. The cost margin
that does exist operates through demand fragmentation---compressed
timelines force smaller orders, mechanically losing bulk
discounts---and through equilibrium shifts in which firm wins each
contract. The reconciliation in Table~\ref{tab:utg_reconciliation}
makes this explicit: of the observed gap, the mechanical bulk-discount
channel over-explains the magnitude; the within-firm component is
zero; the supplier composition residual offsets the bulk-discount
overshoot to deliver the net observed lit-over-admin gap.

\paragraph{Welfare.}
Applied to the \BPwelfareBoundAnnualSpend{} of annual litigated
spending in S\~ao Paulo \citep{cnj2019judicializacao} and a calibration
anchor of \BPwelfareBoundAdmissibility{} of cases admissible to the
administrative channel under expanded committee capacity, the bounded
UTG implies a per-unit-price welfare bound of
\BPwelfareBound{} per year, with the Lee-endpoint range delivering
$[\BPwelfareBoundCIlow, \BPwelfareBoundCIhigh]$. This is a per-unit-price
margin and excludes welfare losses from constrained order sizes (the
mechanical bulk-discount channel itself), reduced bidder participation,
and officials' search time diverted from ordinary procurement. We
report a single defensible figure rather than the stacked ranges
common in the policy literature.

\paragraph{Policy implications.}
Two reforms address the two channels we identify, with first-order
costs the policy maker can plausibly bound.

On the demand side, \emph{expanding the administrative-request
mechanism} eliminates within-urgent sanction exposure for the
admissible share. Since the within-firm pricing component is zero,
the gain from this reform is realized through the demand-aggregation
margin: admin orders are larger, so the bulk-discount channel works in
the buyer's favor. The instrument is administrative capacity (the
scientific-committee infrastructure), not legal reform; the
admissibility share is the binding parameter. We do not endorse a
specific calibration; we show that under any plausible value the
welfare gain scales with bounded UTG times $\$300\,$M.

On the supply side, \emph{framework agreements for repeat-purchase
commodity-like medications} address demand fragmentation directly by
allowing aggregation across urgent and ordinary requests. The instrument
is contract-design reform, available under Lei~14.133/2021. We do not
attach a recovery point estimate to this channel: any number requires
a calibration anchor on the share of items eligible for framework
agreements, and we are not in a position to provide one without
further institutional input.

\paragraph{Interpretation.}
The findings highlight a tension at the heart of right-to-health
enforcement. Judicial mandates intended to guarantee individual access
impose substantial procurement costs on the state, but---in our
data---the mechanism is not a markup story. It is an aggregation and
sourcing story. Institutional designs that preserve the individual
right while restoring the demand-aggregation and supplier-search
margins compromised under compressed timelines---through framework
agreements, expanded administrative-channel capacity, or both---can
deliver welfare gains for litigants and the broader population
simultaneously. The recent Brazilian procurement reform
(Lei~14.133/2021), which expands electronic auctions and introduces
framework agreements, sets up a natural test of this prediction;
future work using post-reform data can quantify the realized recovery.

\paragraph{Boundary.}
Our results are estimated within S\~ao Paulo's BEC platform, on
pharmaceutical procurement in a single state, in 2009-2019. We do not
extrapolate to other states, to non-pharmaceutical procurement, or to
post-reform institutional arrangements. The Lee bounds apply within
the urgent panel and the within-cell variation tilts toward
higher-volume formulary medications (Online Appendix
Table~\ref{tab:both_types_cells}); the welfare figure assumes that
admin admissibility can be expanded without affecting the per-unit-price
gap, an assumption the manuscript does not test.
